\documentclass[11pt,a4paper]{article}
\usepackage[utf8]{inputenc}
\usepackage[T1]{fontenc}
\usepackage[italian]{babel}
\usepackage{geometry}
\usepackage{enumitem}
\usepackage{titlesec}
\usepackage{Alegreya}
\usepackage{AlegreyaSans}

% Page layout
\geometry{margin=2cm}
\pagestyle{plain}

% Typography

% Section formatting
\titleformat{\section}{\Large\sffamily\bfseries}{}{0em}{}[\titlerule]
\titleformat{\subsection}{\large\sffamily\bfseries}{}{0em}{}
\titlespacing{\section}{0pt}{1.5em}{0.8em}
\titlespacing{\subsection}{0pt}{1em}{0.5em}

% List formatting
\setlist[itemize]{leftmargin=*,nosep,before=\vspace{0.2em},after=\vspace{0.2em}}
\setlist[enumerate]{leftmargin=*,nosep,before=\vspace{0.2em},after=\vspace{0.2em}}

% Custom commands for consistent formatting
\newcommand{\cvitem}[4]{%
    \textbf{#1} \hfill \textit{#2}\\
    #3\\
    #4\vspace{0.3em}
}

\newcommand{\cvevent}[3]{%
    \textbf{#1} \hfill \textit{#2}\\
    #3\vspace{0.3em}
}

\begin{document}

% Header
\begin{center}
    {\huge\sffamily\bfseries ALICE CORTEGIANI}\\[0.5em]
    {\large CURRICULUM VITAE ET STUDIORUM}\\[1em]

    alicecortegiani@gmail.com \quad $\bullet$ \quad +39 3385262404\\
    Vocabolo Madonna del Mattone, SNC -- 02040 Mompeo (RI)
\end{center}

\vspace{1em}

Si diploma in Clarinetto e Musica da Camera presso il Conservatorio di Roma e si perfeziona presso l’Accademia Nazionale di Santa Cecilia con A. Carbonare. La sua attività concertistica l’ha condotta a esibirsi in prestigiose sedi quali Teatro Arsenale per La Biennale di Venezia, Cappella Paolina del Palazzo del Quirinale, Auditorium Parco della Musica di Roma, Teatro dell’Opera di Roma e Auditorio Nacional de Musica de Madrid. La sua attitudine è rivolta verso la ricerca artistica musicale attraverso collaborazioni con compositori e laboratori di sperimentazione, interrogando il ruolo dell’interprete come ponte tra tradizione e contemporaneità. Questa ricerca confluisce naturalmente nella sua attività didattica, dove trasmette tanto la padronanza del repertorio classico quanto l’apertura verso altri linguaggi.


\section{Formazione}

\cvitem{Master di II Livello}{Marzo 2022}
{Università degli Studi Roma Tre}
{Teoria e pratica della Formazione alla Musica di insieme (a.a. 2020/2021)\\
Direttore: Prof. Raffaele Pozzi}

\cvitem{Diploma Accademico di II Livello}{5 Maggio 2017}
{Conservatorio ``S. Cecilia'' di Roma}
{Musica da Camera -- 110/110 e Lode\\
Docente: Roberto Galletto}

\cvitem{Diploma di Vecchio Ordinamento}{2 Luglio 2014}
{Conservatorio ``S. Cecilia'' di Roma}
{Clarinetto -- 10/10\\
Docente: Daniele Rossi}

\cvitem{Corso di Perfezionamento ``I Fiati''}{2016--2020}
{Accademia Nazionale di Santa Cecilia, Roma}
{Clarinetto (a.a. 2016/2017, 2017/2018, 2018/2019, 2019/2020)\\
Docente: Alessandro Carbonare}

\cvitem{Programma Erasmus+}{a.a. 2014/2015}
{Real Conservatorio Superior de Música de Madrid}
{Scambio accademico internazionale}

\subsection{Altri Corsi di Perfezionamento}

\cvitem{Accademia di Musica ``Scatola Sonora'', Roma}{a.a. 2017/2018}
{Clarinetto}{Docente: Calogero Palermo}

\cvitem{Accademia del Colibrì, Pescara}{a.a. 2017/2018, 2018/2019}
{Clarinetto}{Docente: Fabrizio Meloni}

\section{Attività Artistiche e Professionali}

\subsection{Concerti}

\cvevent{LAZZARO IN RIPRESA DI VENTI}{22 Settembre 2025}
{Auditorium del Goethe-Institut Rom, Festival ArteScienza 2025\\
Programma: derive su sculture di suono di Mario Bertoncini\\
Arpa Circolare, Die Lyra des Aeolus, Kathedrale, grande Spirale}

\cvevent{Festival ArteScienza 2025}{10 Luglio 2025}
{MUSA Museo degli Strumenti Musicali, Auditorium Parco della Musica, Roma\\
Programma: Luciano Berio\\
Clarinetto soprano in sib}

\cvevent{Festival ArteScienza 2025}{2 Luglio 2025}
{Auditorium del Goethe-Institut Rom, Centro Ricerche Musicali CRM\\
Programma: Artemi-Maria Gioti, BIAS \emph{prima esecuzione italiana}\\
Clarinetto basso, sistemi interattivi}

\cvevent{LAZZARO 290425}{29 Aprile 2025}
{Teatro Studio, Auditorium Parco della Musica, Roma\\
Seminario Concertato\\
Programma: Domenico Guaccero, Michelangelo Lupone, LAZZARO\\
Clarinetto basso, computer, clavette, live camera}

\cvevent{UMANO POST UMANO}{19 Dicembre 2024}
{MACRO TESTACCIO, La Pelanda, Roma\\
61° Festival di Nuova Consonanza\\
Programma: Agostino Di Scipio, Umano Post Umano \emph{prima esecuzione assoluta}\\
Clarinetto basso, computer, clavette, live camera}

\cvevent{LAZZARO101024}{10 Ottobre 2024}
{Auditorium del Conservatorio ``A. Casella'' di L'Aquila\\
ElettroAQustica\\
Programma: Luigi Nono, Michelangelo Lupone, LAZZARO\\
Clarinetto contrabbasso}

\cvevent{LAZZARO130724}{13 Luglio 2024}
{Auditorium del Goethe-Institut Rom -- ArteScienza 2024, Centro Ricerche Musicali CRM\\
Programma: Luigi Nono, Michelangelo Lupone, LAZZARO\\
Clarinetto contrabbasso, clarinetto basso e filtro a pedale}

\cvevent{LAZZARO220324}{22 Marzo 2024}
{Conferenza sulla musica elettronica -- Deutsches Historisches Institut in Rom\\
Programma: Domenico Guaccero, LAZZARO\\
Clarinetto contrabbasso}

%\subsection{Concerti e Festival 2023}

\cvevent{LAZZARO 030923 \& 111023}{3 Settembre / 11 Ottobre 2023}
{Festival ArteScienza 2023 / Festival ``Le Forme del Suono''\\
Auditorium Goethe-Institut Rom / Conservatorio ``O. Respighi'' di Latina\\
Programma: Walter Branchi, Giuseppe Silvi, Domenico Guaccero, LAZZARO}


\cvevent{Orchestra del Teatro Petruzzelli di Bari}{26 Ottobre -- 3 Novembre 2023}
{Fondazione Lirico Sinfonica Petruzzelli\\
Progetto AUS ITALIEN -- Claudio Ambrosini\\
Ruolo: Clarinetto basso}

\cvevent{60° Festival di Nuova Consonanza}{27 Novembre 2023}
{Teatro Studio del Parco della Musica, Roma\\
Ondrej Adámek: ``Ça tourne Ça bloque'', ``Let me tell you a story'', ``Sinuous Voices''\\
PMCE -- Parco della Musica Contemporanea Ensemble}

%\subsection{Concerti e Festival 2022}

\cvevent{Orchestra Nazionale di Santa Cecilia}{8--12 Dicembre 2022}
{Accademia Nazionale di Santa Cecilia, Stagione Sinfonica 2022-2023\\
Haydn, Kodaly, Dvorak -- Direttore: Antonio Pappano}

\cvevent{Orchestra Nazionale di Santa Cecilia}{18--22 Ottobre 2022}
{Inaugurazione Stagione Sinfonica 2022-2023\\
R. Strauss: \textit{Elektra} -- Direttore: Antonio Pappano}

\cvevent{Festival ArteScienza 2020}{2020}
{Giardini del Goethe-Institut Rom\\
DUO ESSENTIA: Sofia Avramidou, Toshio Hosokawa, Jukka Tiensuu, Dai Fujikura}

\cvevent{Centro Ricerche Musicali CRM}{8 Marzo 2019}
{Goethe-Institut Rom\\
Luciano Berio, Giacinto Scelsi, Karlheinz Stockhausen, Alessandra Ravera, Jörg Widmann}

\cvevent{La Biennale di Venezia}{6 Ottobre 2019}
{Teatro Piccolo Arsenale, Venezia\\
Ensemble Novecento dell'Accademia Nazionale di Santa Cecilia\\
Direttore: Matthieu Mantanus}

\subsection{Masterclass e Docenza}

\cvevent{Masterclass in qualità di Docente}{21/09 -- 16/10/2023}
{Conservatorio ``D. Cimarosa'' di Avellino\\
``LUZ, da Descrizione del corpo'' -- VII edizione progetto ``Le Scritture''\\
Curata da Maria Pia Sepe e Giacomo Vitale}

\cvevent{Masterclass in qualità di Docente}{17 Maggio 2021}
{Conservatorio della Svizzera Italiana di Lugano\\
Classe di composizione del M° Nadir Vassena\\
Tecniche e prassi esecutive nel repertorio moderno e contemporaneo}

\subsection{Registrazioni Discografiche}

\cvevent{CD ``BROKEN SHAKE''}{2020}
{DUO ESSENTIA -- Clarinetto/Clarinetto basso e Fisarmonica\\
Alice Cortegiani, clarinetti \quad Samuele Telari, fisarmonica\\
EMA Vinci L\&C Records -- UPC 0650414049562\\
Concerto di presentazione: 21/05/2021, Sala Casella, Accademia Filarmonica Romana}

%\subsection{Altri Eventi Significativi}



\vfill
\begin{flushright}
    \textit{Roma, 20 Settembre 2025}\\[0.5em]
    \rule{4cm}{0.4pt}\\
    Firma
\end{flushright}

\end{document}
